\input{preamble}
\usepackage{cancel}
\usepackage{tikz}
\usetikzlibrary{positioning}
\usetikzlibrary{shapes}
\usepackage{titlesec}
\titleformat{\section}{\Large\scshape\raggedright}{}{0em}{}[]

\title{}
\author{Dane Jeon}
\affiliation{Sogang University,\\
35 Baekbeom-ro, Mapo-gu, Seoul 04107, Republic of Korea}
\emailAdd{danetrouveur@gmail.com}
\abstract{}

\begin{document}

\section{Pairwise Equal Charge Solution in Kerr-Schild Form}
 After failing to express the 2-charge pairwise equal charge solution in the ansatz structure from Wu, we attempt to express the pairwise equal charge solution in a Kerr-Schild-like form. That is, we attempt to express the metric in the form
\[ds^2=d\bar{s}^2+\dfrac{2m}{U}(k_{\mu}\,dx^{\mu })^2\tag{1}\]
where $d\bar{s}^2$ is a AdS$_{4}$ background metric we expect from the massless limit and the 1-form $k_{\mu }$ is linear to the electric gauge potentials.
\newline
\newline
The logic we follow in this section is quite straight forward -- alike how Wu has sought the Cvetic-Youm solutions, I look at the $g= 0$ limit to obtain the appropriate coordinate transformation. 
\newline

\begin{prop}
	Upon the following reparametrisation of the ungauged metric ($g=0$) with charge
	\[{R^2=(r+2ms_1^2)(r+2ms_2^2)}\tag{2}\]
	we arrive at the following form (the new coordinate $R$ is simply written as $r$)
\[
\begin{align*}
	ds^2=&-dt^2+(r^2+a^2\cos ^2\theta )\Big(\dfrac{dr^2}{\mathrm{\Delta }_{r}}+d\theta ^2\Big)+(a^2+r^2)\sin ^2\theta \,d\phi ^2\\
	&+\dfrac{2mU}{r^2+a^2\cos ^2\theta }(dt-a\sin ^2\theta \,d\phi )^2
\end{align*}
	\tag{3}
\]
with
\[
\begin{align*}
	\mathrm{\Delta }_{r}=&(r^2+m^2(s_1^2-s_2^2)^2)\Big[a^2+r^2-2mU\Big]/r^2\\U=&(1+s_1^2+s_2^2)\sqrt{r^2+m^2(s_1^2-s_2^2)^2}-m[s_1^2(1+s_1^2)+s_2^2(1+s_2^2)]	
\end{align*}
	\tag{4}
\]
The metric shows similaries with that of the ungauged, uncharged Wu form:
\[
\begin{align*}
	ds^2=	&-dt^2+(r^2+a^2\cos ^2\theta )\Big(\dfrac{dr^2}{r^2+a^2-2mr}+\,d\theta ^2\Big)+(a^2+r^2)\sin ^2\theta \,d\phi ^2\\&+\dfrac{2mr}{r^2+a^2\cos ^2\theta }(dt-a\sin ^2\theta \,d\phi )^2
\end{align*}
	\tag{5}
\]

\end{prop}
\vspace{2ex}
\begin{prop}
In the case of the gauged case ($g\ne 0$), the result is similar. Upon the same reparametrisation, we arrive at
\[
\begin{align*}
	ds^2=&-\dfrac{(1+g^2r^2)(1-a^2g^2\cos ^2\theta )}{1-a^2g^2}\,dt^2\\&+(r^2+a^2\cos ^2\theta )\Big(\dfrac{dr^2}{\mathrm{\Delta }_{r}}+\dfrac{d\theta ^2}{1-a^2g^2\cos ^2\theta }\Big)\\&+\dfrac{r^2+a^2}{1-a^2g^2}\sin ^2\theta \,d\phi ^2+\dfrac{2mU}{r^2+a^2\cos ^2\theta }\Big(\dfrac{1-a^2g^2\cos ^2\theta }{1-a^2g^2}\,dt-\dfrac{a}{1-a^2g^2}\sin ^2\theta \,d\phi \Big)^2
\end{align*}
	\tag{6}
\]
with
\[
\begin{align*}
	\mathrm{\Delta }_{r}=(r^2+m^2(s_1^2-s_2^2)^2)\Big[(r^2+a^2)(1+g^2r^2)-2mU\Big]/r^2
\end{align*}
	\tag{7}
\]
and the same $U$ defined above. This again shows a resemblance to the gauged but uncharged solution:
\[
\begin{align*}
	ds^2=&-\dfrac{(1+g^2r^2)(1-a^2g^2\cos ^2\theta )}{1-a^2g^2}\,dt^2\\&+(r^2+a^2\cos ^2\theta )\Big(\dfrac{dr^2}{(r^2+a^2)(1+g^2r^2)-2mr}+\dfrac{d\theta ^2}{1-a^2g^2\cos ^2\theta }\Big)\\
	&+\dfrac{r^2+a^2}{1-a^2g^2}\sin ^2\theta \,d\phi ^2+\dfrac{2mr}{r^2+a^2\cos ^2\theta }\Big(\dfrac{1-a^2g^2\cos ^2\theta }{1-a^2g^2}\,dt-\dfrac{a}{1-a^2g^2}\sin ^2\theta \,d\phi \Big)^2
\end{align*}
	\tag{8}
\]
\end{prop}
\vspace{2ex}

\section{Convention Matching}
With the two independent charge solution perfectly fitting into the Wu form, we now attempt to fit the four charge ungauged case into the Wu form. For this to be properly set up, we first need to match conventions. Using the duality transformation 
\[\begin{cases}
	F^{2}=e^{-\varphi _1+\varphi _2-\varphi _3}\star(\tilde{\mathcal{F}}_{2}+\chi_3e^{2\varphi _3}\mathcal{F}^{1})-\chi_1(\tilde{\mathcal{F}}_{3}+\chi_3F^{4})\\
	F^{3}=e^{-\varphi _1-\varphi _2+\varphi _3}\star(\tilde{\mathcal{F}}_{3}+\chi_2e^{2\varphi _2}\mathcal{F}^{1})-\chi_1(\tilde{\mathcal{F}}_{2}+\chi_2F^{4})
\end{cases}\]
\section{Coefficient Determination Algorithm}
At the current stage of the project, there is a need for the determining the coefficients of the gauge-linear terms. In this section, I suggest that the method of achieving this is by preforming a Taylor series expansion in terms of the charges. Here, we consider the following examples (all of which were proven to be in the Wu form). Ultimately, we will want to apply this logic to the four charge ungauged solution.
\begin{itemize}
	\item[(i)] Three charge AdS$_5$ solution
	\item[(ii)] Two charge AdS$_4$ solution
\end{itemize}
\begin{ex}
(Two charge AdS$_4$) We first present the explicit form of the solution and the gauge fields.
\[\begin{align*}
	ds^2=(H_1H_2)^{1/2}\Big[&-dt^2+\mathrm{\mathrm{\Sigma} }\Big(\dfrac{dr^2}{r^2+a^2-2mr}+d\theta ^2\Big)+(r^2+a^2)\sin ^2\theta \,d\phi ^2\\
	&+\dfrac{s_1^2}{(s_1^2-s_2^2)}\Big(\dfrac{1}{1-1/H_1}\Big)A_1^2+\dfrac{s_2^2}{(s_2^2-s_1^2)}\Big(\dfrac{1}{1-1/H_2}\Big)A_2^2\Big]
\end{align*}\]
Meanwhile, the gauge fields are given as
\[A_{i}=\Big(1-\dfrac{1}{H_i}\Big)\Big(\dfrac{c_i}{s_{i}}\,dt-\dfrac{c_1c_2}{s_{i}c_{i}}a\sin ^2\theta \,d\phi \Big)\]
where
\[\begin{align*}
	H_{i}=&1+\dfrac{2mrs_{i}^2}{\mathrm{\Sigma} }&c_{i}=&\sqrt{1+s_{i}^2}
\end{align*}\]
Preforming a Taylor expansion of the terms inside the warp factor with respect to $s_2$ we find (up to $O(s_2^4)$ terms and subtracting the pure AdS terms)

\[\begin{cases}
	O(0):&\dfrac{1}{s_1^2}\Big(1-\dfrac{1}{H_1}\Big)(c_1dt-c_2a\sin ^2\theta\,d\phi )^2\\
	O(s_2^2):&\dfrac{1}{s_1^{4}}\Big(1-\dfrac{1}{H_1}\Big)(c_1\,dt-c_2a\sin ^2\theta \,d\phi )^2-\dfrac{2mr}{\mathrm{\Sigma} s_1^2}(c_2\,dt-c_1a\sin ^2\theta \,d\phi )^2\\
\end{cases}\]

Doing the same for $s_1$, 

\[\begin{cases}
	O(0):&\dfrac{1}{s_2^2}\Big(1-\dfrac{1}{H_2}\Big)(c_2dt-c_1a\sin ^2\theta\,d\phi )^2\\
	O(s_1^2):&\dfrac{1}{s_2^{4}}\Big(1-\dfrac{1}{H_2}\Big)(c_2\,dt-c_1a\sin ^2\theta \,d\phi )^2-\dfrac{2mr}{\mathrm{\Sigma} s_2^2}(c_1\,dt-c_2a\sin ^2\theta \,d\phi )^2\\
\end{cases}\]

\end{ex}
\vspace{2ex}
\begin{ex}
(Three charge AdS$_5$ solution) We first present the explicit form of the solution and the gauge fields.
\[\begin{align*}
	ds^2=(H_1H_2H_3)^{1/3}\Big[&-dt^2+\mathrm{\Sigma} \Big(\dfrac{r^2\,dr^2}{(r^2+a^2)(r^2+b^2)-2mr}+d\theta ^2\Big)+(r^2+a^2)\sin ^2\theta \,d\phi ^2\\
	&+(r^2+b^2)\cos ^2\theta \,d\psi ^2+\dfrac{s_1^2}{(s_1^2-s_2^2)(s_1^2-s_3^2)}\Big(\dfrac{\mathrm{\Sigma} H_1}{2m}\Big)A_1^2\\
	&+\dfrac{s_2^2}{(s_2^2-s_1^2)(s_2^2-s_3^2)}\Big(\dfrac{\mathrm{\Sigma} H_2}{2m}\Big)A_2^2+\dfrac{s_3}{(s_3^2-s_1^2)(s_3^2-s_2^2)}\Big(\dfrac{\mathrm{\Sigma} H_3}{2m}\Big)A_3^2\Big]
\end{align*}\]
Meanwhile, the guage fields are given as 
\[A_{i}=\dfrac{2m}{\mathrm{\Sigma} H_{i}}\Big[\dfrac{s_{i}c_1c_2c_3}{c_{i}}\Big(\dfrac{c_{i}^2}{c_1c_2c_3}dt-a\sin ^2\theta \,d\phi -b\cos ^2\theta \,d\psi \Big)+\dfrac{c_{i}s_1s_2s_3}{s_{i}}\Big(b\sin ^2\theta \,d\phi +a\cos ^2\theta\, d\psi \Big)\Big]\]
where
\[\begin{align*}
	H_{i}=&1+\dfrac{2ms_{i}^2}{\mathrm{\Sigma} }&c_{i}=&\sqrt{1+s_{i}^2}
\end{align*}\]
Preforming a Taylor expansion of the terms inside the warp factor with respect to $(s_2,s_3)$ we find (up to $O(s_2^{4},s_3^{4})$ terms and subtracting the pure AdS terms)

\[\begin{align*}
	O(0):\dfrac{2m}{\mathrm{\Sigma} }\Big[&\dfrac{1}{H_1}\dfrac{1}{s_1^2}\Big(c_1\,dt-\dfrac{c_1c_2c_3}{c_1}(a\sin ^2\theta \,d\phi +b\cos ^2\theta \,d\psi )\Big)^2+\\ &c_2^2(b\sin ^2\theta \,d\phi +a\cos ^2\theta \,d\psi )^2+c_3^2(b\sin ^2\theta \,d\phi +a\cos ^2\theta \,d\psi )^2\Big]
\end{align*}\]

Doing the same for $(s_1,s_3)$,

\[\begin{align*}
O(0):\dfrac{2m}{\mathrm{\Sigma} }\Big[&\dfrac{1}{H_2}\dfrac{1}{s_2^2}\Big(c_2\,dt-\dfrac{c_1c_2c_3}{c_2}(a\sin ^2\theta \,d\phi +b\cos ^2\theta \,d\psi )\Big)^2+\\ &c_1^2(b\sin ^2\theta \,d\phi +a\cos ^2\theta \,d\psi )^2+c_3^2(b\sin ^2\theta \,d\phi +a\cos ^2\theta \,d\psi )^2\Big]
\end{align*}\]

Doing the same for $(s_1,s_2)$,

\[\begin{align*}
O(0):\dfrac{2m}{\mathrm{\Sigma} }\Big[&\dfrac{1}{H_3}\dfrac{1}{s_3^2}\Big(c_3\,dt-\dfrac{c_1c_2c_3}{c_3}(a\sin ^2\theta \,d\phi +b\cos ^2\theta \,d\psi )\Big)^2+\\ &c_1^2(b\sin ^2\theta \,d\phi +a\cos ^2\theta \,d\psi )^2+c_2^2(b\sin ^2\theta \,d\phi +a\cos ^2\theta \,d\psi )^2\Big]
\end{align*}\]

\end{ex}
\vspace{2ex}
\begin{ex}
(Four charge AdS$_4$ solution) 
We first present the explicit form of the solution and the gauge fields,
\[\begin{align*}
	ds^2=\dfrac{W}{\mathrm{\Sigma} }\Big[&-\dfrac{\mathrm{\Sigma} (r^2+a^2\cos ^2\theta -2mr)}{W^2}\Big(dt+\dfrac{2ma\sin ^2\theta (rc_1c_2c_3c_4-(r-2m)s_1s_2s_3s_4)}{r^2+a^2\cos ^2\theta -2mr}\,d\phi \Big)^2\\
	&+\mathrm{\Sigma} \Big(\dfrac{dr^2}{r^2+a^2-2mr}+d\theta ^2+\dfrac{r^2+a^2-2mr}{r^2+a^2\cos ^2\theta -2mr}\sin ^2\theta \,d\phi ^2\Big)\Big]
\end{align*}\]
Meanwhile, the gauge fields are given as
\[\begin{align*}
	A_i=\dfrac{2m}{aW^2}\Big[&(\dfrac{r_1r_2r_3r_4}{r_{i}}+ru^2)[c_is_ia\,dt-(a^2-u^2)(\dfrac{c_1c_2c_3c_4}{c_{i}}s_i-\dfrac{s_1s_2s_3s_4}{s_{i}}c_{i})\,d\phi ]\\
	&+2mu^2[e_ia\,dt-(a^2-u^2)\dfrac{s_1s_2s_3s_4}{s_{i}}c_i\,d\phi ]\Big]
\end{align*}\]
where
	\[e_{i}=\dfrac{c_1c_2c_3c_4s_1s_2s_3s_4}{c_{i}s_{i}}(c_{i}^2+s_{i}^2)-c_{i}s_{i}\Big(\sum _{\substack{j<k\\j,k\ne i}}s_{j}^2s_{k}^2+\dfrac{2s_1^2s_2^2s_3^2s_4^2}{s_{i}^2}\Big)\]
and
\[\begin{align*}
	H_{i}=&1+\dfrac{2mrs_{i}^2}{\mathrm{\Sigma} }
	\quad
	c_{i}=\sqrt{1+s_{i}^2}
	\quad
	r_{i}=r+2ms_{i}^2\\
	W^{2}=&(r_1r_2+u^2)(r_3r_4+u^2)-4mu^2(s_1s_2c_3c_4-c_1c_2s_3s_4)^2
\end{align*}\]
Preforming a Taylor expansion of the terms inside the warp factor with respect to $(s_2,s_3,s_4)$ we find (up to $O(s_2^{4},s_3^{4},s_4^{2})$ terms and subtracting the pure AdS terms)

\[\begin{align*}
	O(0):&-\dfrac{\mathrm{\Sigma} (r^2+a^2\cos ^2\theta -2mr)}{(r_1r+u^2)(r^2+u^2)}\Big(dt+\dfrac{2mrc_1a\sin ^2\theta}{r^2+a^2\cos ^2\theta -2mr}\,d\phi \Big)^2\\
	&+\mathrm{\Sigma} \Big(\dfrac{dr^2}{r^2+a^2-2mr}+d\theta ^2+\dfrac{r^2+a^2-2mr}{r^2+a^2\cos ^2\theta -2mr}\sin ^2\theta \,d\phi ^2\Big)
\end{align*}\]


Doing the same for $(s_1,s_3,s_4)$,

\[\begin{align*}
	O(0):&-\dfrac{\mathrm{\Sigma} (r^2+a^2\cos ^2\theta -2mr)}{(r_2r+u^2)(r^2+u^2)}\Big(dt+\dfrac{2mrc_2a\sin ^2\theta}{r^2+a^2\cos ^2\theta -2mr}\,d\phi \Big)^2\\
	&+\mathrm{\Sigma} \Big(\dfrac{dr^2}{r^2+a^2-2mr}+d\theta ^2+\dfrac{r^2+a^2-2mr}{r^2+a^2\cos ^2\theta -2mr}\sin ^2\theta \,d\phi ^2\Big)
\end{align*}\]

Doing the same for $(s_1,s_2,s_4)$,

\[\begin{align*}
	O(0):&-\dfrac{\mathrm{\Sigma} (r^2+a^2\cos ^2\theta -2mr)}{(r_3r+u^2)(r^2+u^2)}\Big(dt+\dfrac{2mrc_3a\sin ^2\theta}{r^2+a^2\cos ^2\theta -2mr}\,d\phi \Big)^2\\
	&+\mathrm{\Sigma} \Big(\dfrac{dr^2}{r^2+a^2-2mr}+d\theta ^2+\dfrac{r^2+a^2-2mr}{r^2+a^2\cos ^2\theta -2mr}\sin ^2\theta \,d\phi ^2\Big)
\end{align*}\]

Doing the same for $(s_2,s_3,s_4)$,

\[\begin{align*}
	O(0):&-\dfrac{\mathrm{\Sigma} (r^2+a^2\cos ^2\theta -2mr)}{(r_4r+u^2)(r^2+u^2)}\Big(dt+\dfrac{2mrc_4a\sin ^2\theta}{r^2+a^2\cos ^2\theta -2mr}\,d\phi \Big)^2\\
	&+\mathrm{\Sigma} \Big(\dfrac{dr^2}{r^2+a^2-2mr}+d\theta ^2+\dfrac{r^2+a^2-2mr}{r^2+a^2\cos ^2\theta -2mr}\sin ^2\theta \,d\phi ^2\Big)
\end{align*}\]
Simplifying the first expression we found and subtracting the pure AdS part, we find
\[\begin{align*}
	\dfrac{2mr}{\mathrm{\Sigma} H_1}(c_1\,dt-a\sin ^2\theta \,d\phi )^2
\end{align*}\]
\end{ex}
\vspace{2ex}

\section{Ungauged Four Charge Ansatz}
Using the knowledge from the coefficient determination process above and also educated guess work, I was able to arrive at the following form of the ungauged four charge ansatz.
\[\begin{align*}
ds^2=\dfrac{W}{r^2+a^2\cos ^2\theta }\bigg[&-dt^2+(r^2+a^2\cos ^2\theta )\Big(\dfrac{dr^2}{r^2+a^2-2mr}+d\theta ^2\Big)\\&+(r^2+a^2)\sin ^2\theta \,d\phi ^2+\frac{W^2}{2m}\bigg(\sum ^{4}_{i=1}F_{i}A_{i}^2+\sum _{i=1}^{4}(G_{i}-F_{i})[A^{(\phi )}_{i}]^2\bigg)\bigg]
\end{align*}\]
where $A^{(\phi )}_{i}$ denotes the $\phi $-part of the gauge field $A_{i}$ and $F_{i}$ and $G_{i}$ each denote the following:
\[\begin{cases}
F_{i}=\bigg(\dfrac{r_1r_2r_3r_4}{r_{i}}+ra^2\cos ^2\theta +2ma^2\cos ^2\theta\, \dfrac{e_{i}}{c_{i}s_{i}}\bigg)^{-1}s_{i}^{4}\bigg(\displaystyle\prod_{j\ne i}^{}{(s_{i}^2-s_{j}^2)}^{-1}\bigg)\\\\
G_{i}=\Big(\dfrac{r_1r_2r_3r_4}{r_{i}}+ra^2\cos ^2\theta +2ma^2\cos ^2\theta\,\dfrac{s_1s_2s_3s_4c_{i}^2}{c_1c_2c_3c_4s_{i}^2-s_1s_2s_3s_4c_{i}^2}\Big)^{-1}s_{i}^{4}\bigg(\displaystyle\prod_{j\ne i}(s_{i}^2-s_{j}^2)^{-1}\bigg)
\end{cases}\]
\section{General Gauged Four Charge Ansatz}
The ungauged four charge ansatz above compels us to start constructing the gauged four charge ansatz we've been building towards. Due to the elegance of the Wu form, the pure anti-de Sitter part has a trivial gauged generalisation, namely
\[\begin{align*}
ds^2=\dfrac{W}{r^2+a^2\cos ^2\theta }\bigg[&-\dfrac{(1+g^2r^2)(1-a^2g^2\cos ^2\theta )}{1-a^2g^2}\,dt^2+(r^2+a^2\cos ^2\theta )\Big(\dfrac{dr^2}{\mathrm{\Delta}_{r}}+\dfrac{d\theta ^2}{1-a^2g^2\cos ^2\theta }\Big)\\&+\dfrac{r^2+a^2}{1-a^2g^2}\sin ^2\theta \,d\phi ^2+\dfrac{W^2}{2m}\Big(\sum ^{4}
_{i=1}F_{i}A_{i}^2+\sum ^{4}_{i=1}(G_{i}-F_{i})[A_{i}^{(\phi )}]^2\Big)\bigg]
\end{align*}\]
and the new gauge fields would be defined through the equality
\[\begin{align*}
\dfrac{W^2}{2m}A_{i}=&\Big(\dfrac{r_1r_2r_3r_4}{r_{i}}+2ma^2\cos ^2\theta \dfrac{e_{i}}{c_{i}s_{i}}+ra^2\cos ^2\theta \Big)\\&\times \Big[K_1\,c_{i}s_{i}\,dt-K_2\,a\sin^2 \theta \Big(\dfrac{c_1c_2c_3c_4s_{i}}{c_{i}}-\dfrac{s_1s_2s_3s_4c_{i}}{s_{i}}\Big)\,d\phi\Big]\\ &+K_3\,2ma^2\cos ^2\theta \Big(\dfrac{c_1c_2c_3c_4}{c_{i}^2}e_{i}-\dfrac{s_1s_2s_3s_4}{s_{i}^2}(e_{i}+c_{i}s_{i})\Big)a\sin ^2\theta \,d\phi 
\end{align*}\]
Notice that we have a total of 5 undetermined elements, namely $(\mathrm{\Delta }_{r},K_1,K_2,K_3)$. We suggest the following methods for the determination process:
\begin{itemize}
\item[(i)] Proper truncation limits for the two independent charge case and the pairwise identification case
\item[(ii)] Educated guesswork using symmetry and the original Wu paper
\end{itemize}
We will try to employ these techniques in this exact order. 
\vspace{2ex}
\section{Gauged Electric Potentials}
\begin{defi}
(Discrete inversion symmetry) The solution is expected to be invariant under the transformation
\[a\rightarrow \dfrac{a}{(ag)^2}\quad r\rightarrow \dfrac{r}{(ag)}\quad a\cos \theta \rightarrow \dfrac{a\cos \theta }{(ag)}\quad m\rightarrow \dfrac{m}{(ag)^3}\quad \phi \leftrightarrow gt\quad  s_{i}\rightarrow (ag)s_{i}\]
We define the weight of some function $w(f)$ as the power of $ag$ it obtains as it goes through the transformation.
\end{defi}
\vspace{2ex}
\begin{lem}
(Symmetrising $W$) We try to determine the gauged $W$ through matching the weights of the term as they go through the inversion transformation denoted above. Consider the following function
\[\begin{align*}
W^2= r_1r_2r_3r_4+u^{4}+u^2\bigg[&2r^2+2mr(s_1^2+s_2^2+s_3^2+s_4^2)
+8m^2c_1c_2c_3c_4s_1s_2s_3s_4\\&-4m^2(s_{123}^2+s_{124}^2+s_{134}^2+s_{234}^2+2s_{1234}^2)\bigg]
\end{align*}\]
we can make this function uniformly of weight $w(W^2)=-4$ by writing
\[\begin{align*}
W^2=&r_1r_2r_3r_4+u^{4}+u^{2}\bigg[2r^2+2mr(s_1^2+s_2^2+s_3^2+s_4^2)+8m^2c_{1234}s_{1234}\\
&-4m^2s_{1234}\sqrt{(c_1^2-1)(c_2^2-1)(c_3^2-1)(c_4^2-1)}\Big(\dfrac{1}{c_4^2-1}+\dfrac{1}{c_3^2-1}+\dfrac{1}{c_2^2-1}+\dfrac{1}{c_1^2-1}+2\Big)\bigg]
\end{align*}\]
If we additionally impose (naturally, as our metric has a factor $W/(r^2+a^2\cos ^2\theta )$) that $W^2\rightarrow W^2/(ag)^4$, we need to multiply each $c_{i}$ by $\tilde{c}_{i}$, giving 
\[\begin{align*}
W^2=&r_1r_2r_3r_4+u^{4}+u^{2}\bigg[2r^2+2mr(s_1^2+s_2^2+s_3^2+s_4^2)+8m^2\tilde{c}_{1234}c_{1234}s_{1234}\\
&-4m^2s_{1234}\Big(\prod_{i=1}^{4}\sqrt{\tilde{c_i}^2c_i^2-1}\Big)\Big(\dfrac{1}{\tilde{c}_{4}^2c_4^2-1}+\dfrac{1}{\tilde{c}_{3}^2c_3^2-1}+\dfrac{1}{\tilde{c}_{2}^2c_2^2-1}+\dfrac{1}{\tilde{c}_{1}^2c_1^2-1}+2\Big)\bigg]
\end{align*}\]
The uniqueness of this invariant form must be thought of further.
\end{lem}
\vspace{2ex}
\begin{prop}
(Ansatz for symmetric $A_{i}$) We try to determine the gauged $A$ by going through a similar process denoted above. Consider the ungauged electric potential
\[\begin{align*}
A_1=\dfrac{2m}{W^2}\bigg[&(r_1r_2r_3+ra^2\cos ^2\theta )\underbrace{(c_1s_1\,dt-a\sin ^2\theta (c_2c_3c_4s_1-s_2s_3s_4c_1)\,d\phi )}_{(\mathrm{I})}\\
&+2ma^2\cos ^2\theta \underbrace{(e_1\,dt-a\sin ^2\theta s_2s_3s_4c_1\,d\phi )}_{(\mathrm{II})}\bigg]
\end{align*}\]
Firstly, we must notice that a overall factor of $1/\mathrm{\Xi }$ is obvious, given the need of a sign flip under inversion. In the process of employing inversion symmetry, we divide the gauge fields into two parts, (I) and (II). Gauging part (I) is relatively simple. We need to start by first noticing that the overall coeffients simply have a weight of 0, and that we expect the section to map onto itself perfectly up to a sign difference.  
\[\begin{align*}
A_{1}^{(\mathrm{I})}=(c_1s_1\tilde{c_2}\tilde{c_3}\tilde{c_4}-(ag)^2\tilde{c}_1s_2s_3s_4)\mathrm{\Delta }_{\theta }\,dt-a\sin ^2\theta (\tilde{c_1}c_2c_3c_4s_1-s_2s_3s_4c_1)\,d\phi \end{align*}\]
\end{prop}
\vspace{2ex}
Here, a problem is that any factor of $\tilde{c}_{i}$ can add to the last term, which needs further investigation. What about part (II)? This time, the overall factors scale by $-2$ and we are compelled to think that part (II) should scale by 2. We conclude
\[A^{(\mathrm{II})}_{1}=e_1 \mathrm{\Delta }_{\theta }\,dt-\Big[(\tilde{c}_{1}^2+(ag s_1)^2)\tilde{c_2}\tilde{c_3}\tilde{c_4}s_2s_3s_4-(a g)^2\tilde{c_1}s_1(s_2^2s_3^2+s_2^2s_4^2+s_3^2s_4^2+2(ag)^2s_2^2s_3^2s_4^2)\Big]a\sin ^2\theta \,d\phi \]
for $e_1=c_2 c_3 c_4 s_2 s_3 s_4 (c_1^2+s_1^2)-c_1s_2 (s_2^2 s_3^2+s_2^2 s_3^2+s_3^2 s_4^2+2 s_2^2 s_3^2 s_4^2 )$. Notice that again, we have the freedom to add $\tilde{c}_{i}$ terms to any term dependent of the factor $(ag)$. Momentarily ignoring the fact, generalising the two, we find
\[A_{i}^{(\mathrm{I})}=\Big(c_is_i\dfrac{\tilde{c_1}\tilde{c_2}\tilde{c_3}\tilde{c_4}}{\tilde{c}_{i}}-(ag)^2\tilde{c}_{i}\dfrac{s_1s_2s_3s_4}{s_{i}}\Big)\mathrm{\Delta }_{\theta }\,dt-\Big(\tilde{c}_{i}s_{i}\dfrac{c_1c_2c_3c_4}{c_{i}}-c_i\dfrac{s_1s_2s_3s_4}{s_{i}}\Big)a\sin ^2\theta \,d\phi \]
and
\[A_{i}^{(\mathrm{II})}=e_{i}\mathrm{\Delta }_{\theta }\,dt-\Big[(\tilde{c}_{i}^2+(ags_{i})^2)\dfrac{\tilde{c_1}\tilde{c_2}\tilde{c_3}\tilde{c_4}}{\tilde{c_{i}}}\dfrac{s_1s_2s_3s_4}{s_{i}}-(ag)^2\tilde{c}_{i}s_{i}\Big(\sum _{\substack{j<k\\j,k\ne i}}s_{j}^2s_{k}^2+2(ag)^2\dfrac{s_1^2s_2^2s_3^2s_4^2}{s_{i}^2}\Big)\Big]a\sin ^2\theta \,d\phi \]
It is easy to verify that these expressions accurately reduce to the ungauged limit for the four charge case. 
\end{prop}
\vspace{2ex}
\begin{rmk}
At this point, we've found a parity conforming ansatz for our electric potentials $A_{i}$. However, recall that our radial functions $F_{i}$ and $G_{i}-F_{i}$ were also functions of charge which were derived from a particular linearisation process. 
\[\dfrac{W^2}{2m}\Big(\sum ^{4}_{i=1}F_{i}A_{i}^2+\sum ^{4}_{i=1}(G_{i}-F_{i})[A_{i}^{(\phi)}]^2\Big)\]
In order to gauge these functions, we apply to logic to the gauged electric potentials. An ultimate test that this sector of the ansatz must pass would be the fact that it would be invariant under parity. Including this certain parity test, we have the following limits that the we want to use to verify and edit our current ansatz and gauge potentials.
\begin{itemize}
\item[(i)] Correct independent charge limit
\item[(ii)] Correct two pairwise equal charge limit
\item[(iii)] Symmetric under parity
\end{itemize}
Well if our current form or perhaps a bit of an improvement yields these results (especially the last one) correctly, it would be really nice! 
\end{rmk}
\vspace{2ex}
\begin{prop}
(Two independent charge limit) We see that under the two independent charge limit, the gauged electric potential reduces to
\[A_1=\dfrac{2m}{W^2}\Big[(r_1r_2r_3+ra^2\cos ^2\theta )A^{(\mathrm{I})}_{1}+2ma^2\cos ^2\theta A^{(\mathrm{II})}_{1}\Big]\]

\end{prop}
\vspace{2ex}

\end{document}
